% Working draft of discussion
\chapter{Discussion}
\label{ch:discussion}

The studies do not establish that more executable evidence generally improves prediction of programming performance. The external-model benchmark correction
showed that a measurement fault could create an apparently large benefit. The bounded-update study then showed that limiting one observation's influence
could help in a controlled tracker stress test. The probe-selection study showed the remaining difficulty: a cautious reliability model still selected
probes less effectively on the held-out average than uncertainty-only selection and the plug-in clipped-belief comparator.

The proposed selector's poorer performance supports keeping uncertainty-only selection as a baseline when testing more
complex rules. The result also shows why each rule needs a comparison under the conditions where it will be used.
Better reliability information and a different decision rule require further evaluation.

\section{What the Mathematical Review Changes}

Review after the experiment identified a gap between the acquisition score and the claimed expected classification benefit. Clipped
beliefs were substituted into a Bayes-risk expression even though they need not be the posterior of the generating model.
The counterexample in Chapter~\ref{ch:method} shows that the score can be positive when neither possible probe outcome changes
the predicted class. Ranking also evaluates hypothetical reliability-specific updates, while execution uses calibration means.

The experiment therefore evaluates the implemented ranking heuristic, not an exact value-of-information method. This narrows
the interpretation of its failure. It cannot establish that accounting for reliability uncertainty is generally ineffective,
or that a correctly specified expected-loss selector would also lose. Conversely, identifying the score limitation does not
establish that repairing it would improve the result. No such repair was evaluated in the frozen study.

The reference with access to true channel rates is similarly limited by its update and scoring rule. Its performance cannot
be used as a bound on the best achievable accuracy. The measured classification errors remain observations of the frozen
implementations. The report's corrected method descriptions preserve those observations.

\section{Alternative Explanations for the Negative Result}

The score mismatch is one possible contributor to the proposed selector's losses, but the experiment does not identify a single cause.
Uncertainty-only selection has an important advantage in this design: the supplied starting probabilities are calibrated,
so closeness to one half directly identifies cases with high initial classification risk. This gives the simple baseline
a useful signal before any reliability estimate is considered. It does not explain every ranking difference or establish
that the same advantage would persist with poorly calibrated starting probabilities.

The lower-quantile score may also down-rank useful checks when their reliability estimates are uncertain. That is the
intended caution mechanism, but caution is not automatically helpful at an exact allocation. The policy must still select
twenty cases. Rejecting one candidate moves the check to another. The study therefore cannot
separate a potentially useful refusal to probe from an unhelpful redistribution of a compulsory budget.

Clipping introduces another trade-off. It limits movement from a bad observation, but can prevent a useful observation from
changing a prediction. The development ablation varied clipping in both hypothetical scoring and executed updating, so it
does not isolate those two effects. Its small descriptive interaction offers no clear evidence that lower-quantile scoring
and clipping reinforce each other. A causal explanation of their respective contributions would require a separately
specified component study. The existing diagnostics are insufficient.

These alternatives change the interpretation of failure. The result rejects the proposed advantage for this combination
of data-generating rules, calibration fit, score and allocation. It does not isolate reliability uncertainty as the cause,
establish that a different quantile would succeed, or justify selecting a better-looking configuration after the fact.

\section{Contribution in Relation to Existing Methods}

Executable evidence and educational prediction were already connected in TIKTOC, while TRAVER examined verification in
coding-tutor dialogues \parencite{duan2025tiktoc,wang2025traver}. ECED supplies a theoretical treatment of sequential noisy-test
selection under explicit assumptions \parencite{chen2017noisy}. The dissertation neither replaces those methods nor provides
a common-benchmark comparison against them. Its combination of familiar components is not sufficient evidence of novelty.

The contribution is a documented analysis of failures in assessment and evidence selection. Predictions were retained before outcome reveal. Replay
exposed how a scoring fault changed the apparent effect. Controlled comparisons tested a proposed response to unreliable
evidence, and a mathematical check exposed a mismatch between the selector's score and its executed decision. Each step can
be examined separately. The value is the documented mechanism and the limits of the resulting claim, rather than a claim
that the pipeline or every safeguard is unprecedented.

This matters for future tutoring systems because choosing to collect more evidence adds another measurement process whose
quality must be checked. The present work gives concrete examples of that problem. It does not establish its prevalence,
show that these faults dominate deployed systems, or demonstrate that the proposed policy solves it.

\section{Selection, Updating and the Application Context}

At the full probe allocation, ordinary selectors request every case and apply the same bounded update. The harm at this
endpoint concerns the use of misleading evidence, independently of how cases are ranked. At smaller allocations, ranking
changes which cases receive that update. These are distinct questions, and a budget curve alone does not isolate every cause
of error.

The starting probabilities in the probe-selection study are calibrated by construction. Its primary result concerns evidence-channel failures,
not erroneous initial learner estimates. Selection is also performed once for a complete batch, with exactly 20 checks assigned
to each episode. The study does not demonstrate a tutor choosing its next question after each student answer, deciding to stop,
or adapting its teaching strategy over time.
Those capabilities remain application motivations rather than measured contributions.

The log-odds bound limits each observation's contribution. It does not prevent repeated misleading observations from accumulating
harm. The bounded-update study improved the adverse-condition average but worsened missing and contradictory conditions. Both results argue for
reporting condition-level failures alongside an aggregate, rather than interpreting a bounded update as a safety guarantee.

\section{What the Historical Data Add}

The exploratory CSEDM analysis asks whether the simplest selection signal remains useful beyond an authored simulator.
Here, the probabilities come from a logistic model fitted to historical learner activity. The analysis covers 729 predictions
from 86 learners across 19 Python problems.

Within the fixed fold-level allocations, selecting the 361 most uncertain predictions captured 141 of 195 model errors
(72.3\%). Random selection at the same allocation would capture 49.5\% in expectation. This provides evidence that
uncertainty can help locate prediction errors in these records, even though it does not explain how to correct them.

This is the connection to the earlier studies. The external-model benchmark tests the value of additional evidence, while
the simulations examine how evidence is interpreted and selected. The historical records test whether uncertainty identifies
where closer attention might be useful. They contain no matching optional-probe outcomes, so the result cannot validate the
probe-selection rule or establish that the simulator's failure settings reflect classrooms.

No reviewer intervened on the selected predictions. Error capture measures where errors were concentrated, rather than how many
could be corrected. Learner-separated folds assess predictions on the existing problem set. They do not test transfer to unseen
problems or courses. The small descriptive accuracy difference between the logistic model and problem-frequency baseline also
does not establish a meaningful predictive advantage by itself.

\section{Implications for Tutor Design}

An executable result needs an inspectable assessment record. The prompt, submitted code, tests, test output and task requirements
should remain linked so that a teacher or learner can see what was checked. The harness correction shows that repeatable execution
can still produce a misleading outcome when comparison logic or test coverage is wrong. A tutor should distinguish code that
failed to run, code that ran and failed a test, and code that passed incomplete tests. It should also provide a route for disputing
the interpretation.

The prediction update should expose its basis. A reviewer should be able to see the previous estimate, the evidence category,
the assumed reliability and the resulting decision. Bounding one update can limit the effect of a bad observation, but the
simulation shows that it can also discard useful information and that repeated errors can accumulate. Simple comparators such as
uncertainty-only selection should remain part of evaluation whenever a more elaborate policy is introduced.

Programming evidence can contain personal work, identifiers or sensitive context. Sending it to an external model and executing
it in a remote sandbox create separate data and security risks. A deployed system would need clear data destinations, limited
retention, deletion arrangements and isolated execution. The recorded fallback keeps the demonstration available during service
failure. Accepting untrusted public submissions would require a separate security design. These concerns sit alongside the transparency, privacy and
beneficence issues identified in reviews of educational LLM systems \parencite{yan2024practical}.

Repeated checks also use learner time. Their burden may fall unevenly if prediction errors differ across prior experience,
language background, disability or access conditions. Before classroom use, the system would need accessible alternatives,
subgroup error analysis and teacher oversight of disputed or high-consequence decisions. The scores studied here provide no
basis for grading a learner, withholding teaching or deciding that a concept has been mastered.

\section{Changes from the Proposal}

The original proposal joined a Socratic multi-agent tutor, a cognitively constrained student simulator, hidden-state tracking and
reinforcement-learning policy evaluation. That design required credible measurements of learning and cognitive load before
its policy results could be interpreted. The first implementation decision was therefore to establish a working tutoring path
and test whether its evidence could support a reliable prediction. The later tracker, acquisition and historical-data studies

% Address reviewer concern: why did Huber bounds help while greedy selection failed?
